% Generated draft: requires mathematical and source review.
\documentclass[11pt]{article}
\usepackage[margin=25mm]{geometry}
\usepackage{fontspec}
\setmainfont{Latin Modern Roman}
\usepackage{amsmath,amssymb,mathtools,unicode-math}
\setmathfont{Latin Modern Math}
\usepackage{enumitem,needspace,xurl,hyperref}
\setlength{\parindent}{0pt}
\setlength{\parskip}{5pt}
\setlength{\emergencystretch}{3em}
\newcommand{\N}{\mathbb{N}}
\newcommand{\Z}{\mathbb{Z}}
\newcommand{\Q}{\mathbb{Q}}
\newcommand{\R}{\mathbb{R}}
\newcommand{\C}{\mathbb{C}}
\newcommand{\F}{\mathbb{F}}
\newcommand{\A}{\mathbb{A}}
\newcommand{\PP}{\mathbb{P}}
\newcommand{\E}{\mathbb{E}}
\newcommand{\eps}{\varepsilon}
\newcommand{\dd}{\,\mathrm{d}}
\begin{document}
\section*{UnsolvedMath ID: 800428}
\textbf{Source problem code:} catalogue entry
\subsection{Definitions and mathematical statement}
Let \(f\in\Q(x)\) be a rational function of degree at least \(2\), regarded as a map \(f:\PP^1(\Q)\to\PP^1(\Q)\). If \(f=P/Q\) with coprime polynomials \(P,Q\in\Q[x]\), its degree is \(\max\{\deg P,\deg Q\}\). Write \(f^0\) for the identity and \(f^{n+1}=f\circ f^n\). A point \(\gamma\in\PP^1(\Q)\) is \emph{\(f\)-wandering} if its forward orbit \(\{f^n(\gamma):n\geq 0\}\) is infinite.

\textbf{Conjecture 51, as supplied.} For every rational function \(f\in\Q(x)\) of degree at least \(2\) and every pair of \(f\)-wandering points \(\alpha,\beta\in\Q\), the Zsigmondy set of \(\mathcal A(f,\alpha,\beta)\) is finite:
\[
\#\operatorname{Zsig}\bigl(\mathcal A(f,\alpha,\beta)\bigr)<\infty.
\]
\textbf{Unresolved definition.} The supplied record does not define \(\mathcal A(f,\alpha,\beta)\), its indexing, or the primitive-divisor convention used to form its Zsigmondy set. For a sequence \((a_n)_{n\geq 1}\) of nonzero integers, the usual convention calls a prime \(p\) a primitive prime divisor of \(a_n\) if \(p\mid a_n\) and \(p\nmid a_m\) for every \(1\leq m<n\); the corresponding Zsigmondy set consists of indices whose terms have no primitive prime divisor. This explains the terminology only and does not identify the unspecified object \(\mathcal A(f,\alpha,\beta)\) with such a sequence. Its construction and conventions must be recovered from the cited source before the conjecture has a fully self-contained formulation.
\subsection{Short English statement}
The supplied conjecture says that every rational map of degree at least two and every pair of rational wandering points give an object \(\mathcal A(f,\alpha,\beta)\) with finite Zsigmondy set. The record omits the definition of this object and its primitive-divisor conventions, so the precise arithmetic assertion remains incomplete.
\subsection{Sources}
\begin{itemize}
\item Local UnsolvedMath registry: permanent ID \(800428\), source problem code \texttt{catalogue entry}, title \emph{Open Problems in Arithmetic Dynamics --- Conjecture 51}. The supplied registry labels the problem \emph{open}; this is a supplied status, not an independently verified literature assessment.
\item Benedetto, DeMarco, Ingram, Jones et al., \emph{Current Trends and Open Problems in Arithmetic Dynamics} (2018), Conjecture 51. Attribution and proposed year are supplied by the registry. Source URL: \url{https://arxiv.org/abs/1806.04980}.
\item Supplied provenance records access on \texttt{2026-10-16} and extraction by \texttt{arxiv\_latex}. These are supplied provenance statements; the source has not been independently accessed or verified for this draft.
\end{itemize}
\end{document}
